% !TeX program = xelatex
% ==========================================================================
% Пропоузал семестрового проєкту — курс «Штучний інтелект», УКУ.
% Це ЄДИНИЙ файл, який вам потрібно редагувати. Уся верстка й брендинг
% живуть у ucuproposal.cls — не змінюйте шрифти, поля, кеглі й кольори.
%
% Компіляція: XeLaTeX + Biber (Overleaf робить це автоматично).
% Локально:   latexmk (див. README.md).
% ==========================================================================
\documentclass[ukrainian]{ucuproposal}
% Опції класу (можна комбінувати через кому):
%   ukrainian | english  — мова документа й логотипу (за замовчуванням ukrainian)
%   slogan                — додає гасло УКУ в колонтитул титулки (за замовчуванням вимкнено)

% --------------------------------------------------------------------------
% Метадані титульної сторінки
% --------------------------------------------------------------------------
\projecttitle{Розробка системи NER моделей для історичних досліджень}
\projecttagline{Аналіз текстів на імена в документах Львова 15-17 століття}

\teammember{Депутат Антон}
\teammember{Головін Максим}
% \teammember{Прізвище Ім'я} % третій учасник (опційно)
% \teammember[https://linkedin.com/in/...]{Прізвище Ім'я} % опційне клікабельне посилання на профіль (LinkedIn тощо)

\mentor{Пахолок Віктор} % якщо ментора немає — видаліть текст і залиште \mentor{}, рядок не друкується

\track{research} % research | engineering — не друкується на титулці, лише для внутрішньої категоризації

\repo{https://github.com/Lloydwqe23/ner\_history\_analysis}

\proposaldate{2026-09-15}

\begin{document}
\maketitlepage

\section{Проблема та мотивація}
% Що саме ви розв'язуєте і чому це важливо й цікаво в контексті ШІ?
% Орієнтовно 1 абзац.
Наш проєкт вирішує проблему визначення персон з даних 15–17 століття у Львові для подальшого вивчення історичної взаємопов'язаності людей у той період. Це важливо для багатьох істориків через велике різноманіття національностей, яке перебувало у Львові на той час. Наш проєкт використовує NER (Named Entity Recognition) моделі для правильного визначення людей у документах. В контексті ШІ цей проєкт допомагає розглянути використання моделей для гуманітарних цілей і об'єднати їх з історичними дослідженнями Львова, і саме цим цей проєкт нас зацікавив. Наша ж мотивація полягає з загальним інтересом в роботі з історично направленим проектом.

\section{Огляд джерел та наявних рішень}
% Що вже зроблено у світі за цією темою? Спирайтесь на references.bib,
% цитуйте через \cite{}. Орієнтовно 3-5 джерел.

Ми здійснили пошук різних латинських датасетів та моделей натренованих на них, таких як efontes/efontes-ner та magistermilitum/roberta-multilingual-medieval-ner. Ми провели оцінку цих моделей за їхньою F1 і швидкістю на тестах, створених за наданими нам текстовими даними, і порівняли рівень помилок цих моделей за F1 і швидкістю.~\cite{efontes_ner,magistermilitum_roberta_medieval,alcar_ner, medieval_multilingual_ner, lasla_latin, ibm_ner}.

\section{Дані}
% Які набори даних, звідки, ліцензія. Якщо збираєте самі — як саме.
% Якщо дані не потрібні — поясніть чому.
Дані взяті з львівських документів ради в періоді з 15 до 17 століття, надані нам із Центру цифрової гуманітаристики. Дані в закритому доступі, за потреби звертатись в центр. В даний момент інша команда займається прочитанням великої кількості сканів, але на момент написання цього пропоузалу обсяг прочитаних і перевірених сягає 150 сторінок, а кількість анотованих (тобто перетворених у чистий текст і вже з виділеними іменами) обмежений 60 сторінками. При потребі ми можемо використати записи Буллінгера~\cite{bullinger_digital} — датасет, що має схожий період написання і схожу версію латині, через що буде хорошим додатком при нестачі матеріалу для тренування NER-моделі. 

\section{Підхід і методи}
% Які алгоритми/архітектури плануєте використати, чи спираєтесь на
% існуючі реалізації та як плануєте їх покращити.
Існують дві основні архітектури NER-моделей: трансформерна і словникова. Трансформери розпізнають сутності завдяки механізму уваги, який аналізує роль кожного слова відносно контексту речення. Словниковий метод полягає в зіставленні частин тексту з попередньо сформованими словниками сутностей за правилами пошуку збігів. Ми постараємось використати найкраще з обох світів, вибравши кращих представників з кожної архітектури, і потім, використовуючи платформу Google Colab і Kaggle, будемо тренувати NER-моделі для покращення їх роботи з нашим датасетом.

\section{Оцінювання результатів}
% Метрики, baseline, спосіб порівняння, якісні очікування.
Для об'єктивної оцінки ефективності ми будемо поділяти розмічені дані на навчальну, валідаційну та тестову вибірки, а як метрику будемо використовувати F1, що є гармонійним середнім між precision і recall. Наше завдання — значно підняти(5-10 відсотків) F1 відносно 80 відсотків, які видають нам стандартні NER-моделі, не дотреновані на наших даних. Ми також будемо використовувати швидкість як метрику, хоч і з меншим пріоритетом, ніж F1, через високу пріоритезацію точності в історичних контекстах (історики можуть почекати день, щоб програма стовідсотково анотувала текст, але вони не можуть покладатись на програму з 40 відсотками точності). Будуть також розглядатися edge cases, де різні моделі постійно помиляються, і ми старатимемось об'єднати їх так, щоб закрити максимальну кількість слабких сторін.

\section{План і ризики}
% Етапи до midterm і до фінального захисту, ключові ризики та план Б.
План складається з 4 пунктів: 1) Ми оберемо декілька моделей, які були натреновані на схожих датасетах (ми на цей момент уже знайшли декілька таких)~\cite{efontes_ner,magistermilitum_roberta_medieval}. 2) Оцінимо F1 і швидкість на тестах та вивчимо глибинно, як працює кожна модель.(вже готова пробна програма для оцінки, яку ми використовували для первинного вибору моделей) 3) Виберемо одну або декілька моделей (залежно від того, чи буде одна добре справлятись), щоб натренувати на нашому датасеті і перевірити працездатність під час роботи з новими даними. 4) Якщо всі попередні пункти будуть виконані, то можна провести дослідження щодо способів утворення зв'язків між людьми за тим, із ким вони зустрічаються в книгах найчастіше, або дослідити, наскільки оптимально працює модель (чи моделі) залежно від кількості даних для тренування. Фінальний результат - це код, що буде використовувати одну чи декілька успішно дотренованих NER-моделей для визначення людей в історичних текстах.

% Приклад врізки з ключовими цифрами (опційно, можна видалити):
% \begin{keyfacts}
% \textbf{Ключові цифри:} 3 датасети \textperiodcentered\ 2 baseline-моделі \textperiodcentered\ 4 тижні до midterm
% \end{keyfacts}

\section{Contribution statement}
% Хто за що відповідає. Відсотки не обов'язкові — якщо хочете їх вказати,
% додайте в квадратних дужках: \contribution[50\%]{Прізвище Ім'я}{...}
\begin{contributions}
  \contribution[50\%]{Депутат Антон}{пошук даних, тренування NER моделей, документація}
  \contribution[50\%]{Головін Максим}{пошук NER-моделей та їх тестування}
\end{contributions}

% Примітка: декларація використання AI-інструментів (AI usage) до цього
% пропоузалу не входить — вона живе в README.md вашого репозиторію.

\ucuprintbibliography

\end{document}
